\section{Transport of the finite layers and the canonical unit}
\label{sec:exponent-transport}

The upper null map supplies a retract built from $m$ local $E_h$-module
layers. We now transport those individual layers by the exact functor
on ordinary \(p\)-local spectra
\[
 T:\Sp\longrightarrow\mathcal D_H,\qquad
 T(V)=L_{K(t)}(E_t\wedge V),
\]
with \(H\) acting through the coherent action on \(E_t\).
Applying \(T\) to a local layer means applying it to the underlying
ordinary spectrum of that object of \(\mathcal C_h\). The mixed trace
will make each resulting layer have $H$-exponent at most one.

\begin{proposition}[The cross-height upper exponent]\label{wup:transfer}
Let \(p\geq5\) be prime, with \(t=p-1\) and \(h=p\). One has
\[
 \expn_H\bigl(T(B_h)\bigr)\leq m=p^2+2.
\]
\end{proposition}

\begin{proof}
The inclusion \(\mathcal C_h\to\Sp\) is exact and lax symmetric
monoidal. Since \(E_h\) is local, the localization unit identifies
the included local algebra \(E_h\) with its given ordinary algebra.
Each local layer
\(E_h\otimes_h I_h^{\otimes_h i}\) in
Proposition~\ref{wup:upper-null} consequently has an ordinary
\(E_h\)-module structure. Ordinary smash with \(E_t\) makes it
a \(D=E_t\wedge E_h\)-module with compatible \(H\)-action.

The \(K(t)\)-acyclic spectra form a smash ideal. Thus the
localization into \(\mathcal C_t\) is symmetric monoidal for its
localized tensor, and the inclusion back into ordinary spectra
is lax symmetric monoidal. Applying localization to the algebra
and this module gives, after that inclusion, a \(D'\)-module
structure on the transformed layer in \(\mathcal D_H\).
Proposition~\ref{wup:completed-trace} gives every such layer
\(H\)-exponent at most one.

The functor \(T\) preserves cofiber sequences and retracts.
Apply it to the finite filtration and retraction of
Proposition~\ref{wup:upper-null}, and add the \(m\) layer bounds
using Lemma~\ref{lem:exponent-calculus}. This yields the stated
bound. The construction used ordinary module structures on the
individual layers; it required no preservation by \(T\) of
the \(\mathcal C_h\) tensor product.
\end{proof}

\begin{corollary}[The exponent contradiction for the canonical unit]
\label{wup:contradiction}
Let \(p\geq5\) be prime, put \(t=p-1\) and \(h=p\), and let
\(X=S^\wedge_p\) be the derived \(p\)-completed sphere. The canonical map
\[
 i_{h,X}=L_t(\eta_X^{(h)}):
 L_tX\longrightarrow L_tL_{K(h)}X
\]
has no homotopy retraction in ordinary spectra.
\end{corollary}

\begin{proof}
The functor \(T\) inverts \(K(t)\)-equivalences: the cofiber of
such a map is \(K(t)\)-acyclic, and remains so after ordinary
smash with \(E_t\). Every \(L_t\)-localization unit is a
\(K(t)\)-equivalence. Lemma~\ref{lem:completion} says that the
actual completion \(c:S\to X\) is a \(K(a)\)-equivalence for
each positive \(a\), in particular for \(a=t,h\).
Write $\lambda_Z:Z\to L_tZ$ for the lower localization unit.
Consider the square of actual localization and completion maps
\[
\begin{CD}
 S @>{\eta_S^{(h)}}>> B_h\\
 @V{\lambda_X\,c}VV
 @VV{\lambda_{L_{K(h)}X}\,L_{K(h)}(c)}V\\
 L_tX @>{i_{h,X}}>> L_tL_{K(h)}X .
\end{CD}
\]
It commutes by naturality of the two localization units.
Both vertical arrows become equivalences under \(T\): for the
left one use the \(K(t)\)-equivalences just stated, and for the
right one use that \(L_{K(h)}(c)\) is an equivalence followed
by a \(K(t)\)-equivalence. Moreover \(T(S)\simeq E_t\), with
coherent \(H\)-action, since \(E_t\) is \(K(t)\)-local.

A homotopy retraction of \(i_{h,X}\) would therefore, after
applying \(T\) and these identifications, exhibit \(E_t\) as a
retract of \(T(B_h)\) in \(\mathcal D_H\).
Proposition~\ref{wup:transfer} and preservation of exponent
bounds under retracts would give
\(\expn_H(E_t)\leq m\). The finite-page lower result
Proposition~\ref{prop:lower-exponent} gives the strict inequality
\(\expn_H(E_t)>m\), a contradiction.
\end{proof}

\begin{proof}[Proof of Theorem~\ref{thm:exponents}]
Proposition~\ref{prop:lower-exponent} gives the strict lower bound and
Proposition~\ref{wup:transfer} gives the upper bound, both for
$m=p^2+2$ in the single category $\mathcal D_H$. The preceding corollary
proves that a retraction of the canonical map would turn the lower object
into an $H$-equivariant retract of the upper object and contradict those
bounds. This proves every assertion of the theorem. The lower bound uses
finite-page survival only; the upper bound uses the actual null map and
its finite layers.
\end{proof}
